        	 	\begin{proof}
        	 	
        	 		Case I: We can follow the arguments exactly as in \cite[Cor.~2.6]{BS17} to conclude all the statements of the theorem. \\
        	 		
        	 		
        	 		Case II: 
        	 		In these Fermi-coordinates, we are now in a position to apply Lemma \ref{lemma:local} and it remains to check that the limit $P$ is in fact a non-trivial (namely: it is not flat) and that the final statement of the theorem, concerning the point-scale sequence $(\widehat{p}_k,\widehat{r}_k)$, also holds. 
        	 		
        	 		
        	 		Concerning the first assertion, observe that by Lemma \ref{lemma:local} it suffices to check that $P$ is not stable. 
        	 		We have ${\jepPubScript{Y}}=\varnothing$ since $\widetilde{P}_k$ is stable on all balls of radius $1/4$ (by virtue of our assumption), hence $\widetilde{P}_k$ converges smoothly and graphically on every compact subset of $\Pi_1$ to a connected minimal hypersurface $P$ and if $\partial P\neq\varnothing$ then $P$ meets $\Pi$ orthogonally. Furthermore either $P$ is properly embedded or $P=\Pi$.  
        	 		As the ambient space is simply connected, we can always conclude that $P$ is two-sided and that this convergence happens with multiplicity one.     		
        	 	
        	 	
        	 	Towards a contradiction, suppose that either $P=\Pi$ or
        	 	$$\lambda_1(P\cap B^{\R^{n+1}}_{3/2}(0))\geqslant 0,$$ so that in particular $P$ is strictly stable on $B^{\R^{n+1}}_{4/3}(0)$. Notice that if $P=\Pi$ then the above is trivially true for variations in the whole of $\R^{n+1}$ (i.e. not just in $\Pi_1$). In either case, by smooth multiplicity one convergence we have that $\widetilde{P}_k$ is strictly stable on $B^{\R^{n+1}}_{4/3}(0)$ which implies (scaling back and using the second part of Remark \ref{rmk:coord}) that $P_k$ is stable on $B_{r_k + r_k^2}(p_k)$, a contradiction. In particular $P$ cannot be planar, and we have 
        	 	$$\lambda_1(P\cap B^{\R^{n+1}}_{3/2}(0))\leqslant -2\lambda_{\ast} < 0$$ for some $\lambda_{\ast}>0$, and hence for all $k$ sufficiently large
        	 	
        	 		$$\lambda_1(\widetilde{P}_k\cap B^{\R^{n+1}}_{3/2}(0))\leqslant -\lambda_{\ast}<0.$$ 
        	 		Hence a rescaling argument implies $$\lambda_1(P_k \cap (B_{2r_k}(p_k)))\longrightarrow -\infty.$$ 
        	 		
        	 		Obviously, the fact that for $n=2$ we must obtain a half-bubble relies on the half-space theorem \cite{HM90} (i.e. there cannot be a non-trivial bubble contained in an open half-space.)
        	 		
        	 		Finally, the argument for the last part is similar to the above except we always choose normal coordinates to carry out our blow-ups using Remark \ref{rmk:coord}. At the $r_k$ scale we have that $B_{\widehat{r}_k}(\widehat{p}_k)$ is similar to the balls 
        	 		$$B^{\R^{n+1}}_{\beta_k^{\pm 1}\widetilde{r}_k}(\widetilde{p}_k)$$ with $1\leqslant \widetilde{r}_k \leqslant C$ and $\widetilde{p}_k\in B^{\R^{n+1}}_C(0)$, at which point we can conclude that $\widetilde{r}_k \to \widetilde{r}\in [1,C]$ and $\widetilde{p}_k \to \widetilde{p}\in B^{\R^{n+1}}_C(0)$.    
        	 		Assuming that $(B^{\R^{n+1}}_{2\widetilde{r}}(\widetilde{p})\smallsetminus B^{\R^{n+1}}_2(0))\cap P$ is stable (or empty), we arrive at a contradiction in a similar fashion as above: first this implies that $B_{3\widetilde{r}/2}^{\R^{n+1}}(\widetilde{p})\smallsetminus B_2^{\R^{n+1}}(0)\cap P$ is strictly stable (or empty). But this domain contains the blown-up initial domain on which we are assuming that $P_k$ is unstable. Thus this must be non-empty, and by the smooth multiplicity one convergence we would obtain that $(B_{3\widetilde{r}/2}^{\R^{n+1}}(\widetilde{p})\smallsetminus B_2^{\R^{n+1}}(0))\cap \widetilde{P}_k$ is strictly stable, a contradiction. 
        	 		
        	 		Thus, we have proved that $$\lambda_1 ((B^{\R^{n+1}}_{3\widetilde{r}/2}(\widetilde{p})\smallsetminus B^{\R^{n+1}}_2(0))\cap P)\leqslant -2\lambda_{\ast}<0 $$ as before, for some $\lambda_{\ast}>0$. A rescaling argument concludes the final statement of the Theorem, once again appealing to Remark \ref{rmk:equivballs}.         
        	 \end{proof}
